\documentclass{article}
\usepackage[utf8]{inputenc}
\usepackage{a4wide}

\begin{document}

\section*{เส้นทางสูงต่ำ}

กำหนดพื้นที่เป็นตารางขนาด\texttt{M\ x\ N}ที่ค่าในนั้นเป็นเลขจำนวนเต็ม

ให้หาเส้นทางที่ทำให้ได้คะแนน\ul{\textbf{สูงสุด}}ในการเดินทางจากตำแหน่ง\texttt{(0,\ 0)}ไปยัง\texttt{(M\ –\ 1,\ N\ –\ 1)}โดยสามารถเดินได้สี่ทางทิศเท่านั้น\textbf{(ทแยงไม่ได้)}ทั้งนี้

คะแนนของแต่ละเส้นทางคือ\ul{\textbf{ค่าที่ต่ำที่สุด}}บนเส้นทางนั้น เช่น

คะแนนของเว้นทาง\texttt{8\ →\ 4\ →\ 5\ →\ 9}คือ\texttt{4}\\



\hfill\break



{\def\LTcaptype{none} % do not increment counter

\begin{longtable}[]{@{}

  >{\raggedright\arraybackslash}p{(\linewidth - 4\tabcolsep) * \real{0.3333}}

  >{\raggedright\arraybackslash}p{(\linewidth - 4\tabcolsep) * \real{0.3333}}

  >{\raggedright\arraybackslash}p{(\linewidth - 4\tabcolsep) * \real{0.3333}}@{}}

\toprule\noalign{}

\begin{minipage}[b]{\linewidth}\raggedright

ข้อมูลเข้า

\end{minipage} & \begin{minipage}[b]{\linewidth}\raggedright

ข้อมูลออก

\end{minipage} & \begin{minipage}[b]{\linewidth}\raggedright

คำอธิบาย

\end{minipage} \\

\midrule\noalign{}

\endhead

\bottomrule\noalign{}

\endlastfoot

\begin{minipage}[t]{\linewidth}\raggedright

3 3\\

5 4 5\\

1 2 6\\

7 4 6\\

\strut

\end{minipage} & 4 & \begin{minipage}[t]{\linewidth}\raggedright

\includegraphics[width=1.25in,height=1.25in,alt={Screenshot 2567-04-01 at 02.20.18.png}]{/public/upload/c2d2565838.png}\\

\strut

\end{minipage} \\

\begin{minipage}[t]{\linewidth}\raggedright

2 6\\

2 2 1 2 2 2\\

1 2 2 2 1 2\\

\strut

\end{minipage} & 2 &

\includegraphics[width=2.08333in,height=0.72917in,alt={Screenshot 2567-04-01 at 02.20.42.png}]{/public/upload/cbe665193a.png} \\

\begin{minipage}[t]{\linewidth}\raggedright

6 5\\

3 4 6 3 4\\

0 2 1 1 7\\

8 8 3 2 7\\

3 2 4 9 8\\

4 1 2 0 0\\

4 6 5 4 3\\

\strut

\end{minipage} & 3 & \begin{minipage}[t]{\linewidth}\raggedright

\includegraphics[width=1.47917in,height=1.77083in,alt={Screenshot 2567-04-01 at 02.22.07.png}]{/public/upload/79b1e75e43.png}\\

\strut

\end{minipage} \\

\end{longtable}

}



\subsection*{Input}
บรรทัดแรก เป็นขนาดพื้นที่\texttt{M,\ N}เป็นจำนวนเต็ม

ขนาดระหว่าง\texttt{2\ x\ 2}ถึง\texttt{100\ x\ 100}



บรรทัดถัดไป M บรรทัดเป็น ค่าที่ตำแหน่งต่างๆ ของพื้นที่

โดยมีค่าระหว่าง\texttt{1}ถึง\texttt{99}



\subsection*{Output}
คะแนน\ul{\textbf{สูงสุด}}ในการเดินทาง\\



\end{document}
